\section{SAM3: Segment Anything Model 3}

\subsection{Evolución de SAM3 respecto a SAM1/SAM2}

La serie SAM (Segment Anything Model) es un modelo de segmentación de uso general presentado por Meta. La evolución clave de SAM3 respecto a las dos generaciones anteriores es \textbf{la incorporación de entrada mediante prompt de texto}:

\begin{table}[H]
\centering
\begin{tabular}{lccc}
\toprule
\textbf{Capacidad} & \textbf{SAM1} & \textbf{SAM2} & \textbf{SAM3} \\
\midrule
Segmentación por prompt de punto/caja & \checkmark & \checkmark & \checkmark \\
Segmentación de video & $\times$ & \checkmark & \checkmark \\
Segmentación por prompt de texto & $\times$ & $\times$ & \checkmark \\
Reconocimiento (salida de Label) & $\times$ & $\times$ & \checkmark \\
\bottomrule
\end{tabular}
\caption{Comparación de capacidades de la serie SAM}
\end{table}

\begin{warningbox}{SAM3 no es un VLM}
Aunque SAM3 admite entrada de texto, solo acepta \textbf{frases nominales simples} (como "cat", "person", "shoe"), y no oraciones en lenguaje natural. No admite entradas en forma de lista (como "cat, remote"), ni tampoco modificadores complejos. Cuando es necesario reconocer varias clases de objetos a la vez, hay que recorrer los nombres de las categorías uno por uno.
\end{warningbox}

\subsection{Flujo de uso de SAM3}

El uso básico de SAM3 es muy sencillo:

\begin{enumerate}
    \item Cargar el modelo SAM3 mediante \texttt{build\_model} (aproximadamente 3.3 GB, se descarga automáticamente al directorio de HuggingFace Hub)
    \item Colocar el modelo en el dispositivo CUDA
    \item Construir el Processor, pasándole la imagen y el prompt de texto
    \item Obtener la salida: Mask (máscara de segmentación), Bounding Box y puntaje de confianza
\end{enumerate}

\begin{lstlisting}[caption={Ejemplo básico de uso de SAM3 (pseudocódigo)}]
# Construir el modelo
model = build_sam3_model()
model.to("cuda")

# Construir el Processor
processor = SAM3Processor(model)

# Leer la imagen y hacer inferencia
image = load_image(url)
with torch.inference_mode():
    outputs = processor(image, text_prompt="remote")
    # outputs contiene: masks, bounding_boxes, confidence_scores
\end{lstlisting}

\subsection{Ejemplo de segmentación a nivel de instancia con SAM3}

Tomemos como ejemplo una imagen clásica del conjunto de datos COCO (dos gatos, dos controles remotos):

\begin{itemize}
    \item Entrada \texttt{"remote"}: identifica con precisión los dos controles remotos, cada uno con su propia Mask y numeración (Instance 0, Instance 1), con una confianza de aproximadamente 0.96
    \item Entrada \texttt{"cat"}: identifica con precisión los dos gatos, con una confianza igualmente alta
    \item Detección de varias clases: al recorrer la lista de categorías \texttt{{[}"remote", "cat"{]}} y combinar los resultados, se obtienen las cuatro instancias en total
\end{itemize}

Para escenas más complejas (como varios niños, zapatos, puertas, etc.), SAM3 también logra la segmentación a nivel de instancia, aunque en objetivos complejos (como una puerta parcialmente oculta) la confianza disminuye en cierta medida.

\subsection{Resumen de la sección}

Al incorporar el prompt de texto, SAM3 logra funciones de reconocimiento sin dejar de mantener una capacidad de segmentación de alta calidad, pero su comprensión de texto se limita a frases nominales simples. Para consultas complejas que requieren razonamiento espacial o comprensión semántica, es necesario combinarlo con un VLM para compensar esa limitación.


